\section{Planejamento de trajetória e autonomia}

O planejamento de trajetória em operações espaciais constitui uma etapa central para garantir que a espaçonave, ou o manipulador acoplado a ela, execute a missão de forma segura e compatível com as restrições do problema. Em cenários de serviço em órbita, inspeção, captura e berthing, esse planejamento não se restringe à definição geométrica de um caminho, mas passa a envolver simultaneamente restrições cinodinâmicas, limites de atuação, risco de colisão, incertezas associadas ao alvo e requisitos crescentes de autonomia \cite{santos2022machine,capolupo2017heuristic,fiori2024modeling}.

Entre as abordagens encontradas na literatura, destacam-se as formulações baseadas em controle ótimo e otimização multiobjetivo. No trabalho de \cite{santos2022machine}, o problema de planejamento automático de trajetória de uma espaçonave com manipulador em órbita é formulado de modo a considerar, de forma simultânea, o deslocamento da plataforma, a manipulabilidade do braço, o torque máximo, o posicionamento do efetuador final, a prevenção de colisão entre o manipulador e a espaçonave e a minimização dos esforços de atuação sob incerteza na região de berthing. Essa formulação evidencia que, em robótica espacial, o planejamento de trajetória está fortemente associado ao tratamento simultâneo de múltiplos objetivos de desempenho e de restrições físicas do sistema \cite{santos2022machine}.

Outra classe importante de métodos corresponde às técnicas de motion planning baseadas em amostragem. Em \cite{capolupo2017heuristic}, algoritmos como Rapidly Exploring Random Trees e Sampling Based Model Predictive Optimization são empregados para construir sequências viáveis de manobras em operações de proximidade ao redor de pequenos corpos celestes, respeitando restrições cinodinâmicas e de trajetória. Os autores ressaltam que esse tipo de abordagem é particularmente útil em missões de maior nível de abstração, nas quais o problema não é definido apenas por waypoints previamente fixados, mas por objetivos mais amplos, como observação, mapeamento ou pouso autônomo em ambientes complexos \cite{capolupo2017heuristic}.

A incorporação de aprendizado de máquina ao problema também tem recebido atenção recente. Em \cite{santos2022machine}, uma estratégia de machine learning é empregada para apoiar a análise de cinemática inversa dentro do processo de planejamento ótimo, utilizando dados de treinamento e tarefas anteriores para melhorar a convergência do procedimento de otimização. De modo semelhante, \cite{santos2023physics} propõe uma abordagem physics informed machine learning em que uma Support Vector Machine é combinada a um modelo físico simplificado para estimar boas soluções iniciais para o problema inverso, incorporando efeitos dinâmicos do satélite e do manipulador ao processo de aprendizagem. Em ambos os casos, observa-se uma tendência de integrar modelos físicos e técnicas orientadas por dados para aumentar a eficiência computacional e apoiar níveis mais elevados de autonomia \cite{santos2022machine,santos2023physics}.

Além do planejamento propriamente dito, a literatura também apresenta contribuições relevantes no campo do guiamento autônomo sob restrições posicionais. O trabalho de \cite{fiori2024modeling} propõe um arcabouço de modelagem, simulação e controle para o movimento de uma pequena espaçonave nas proximidades de uma estação espacial, tratando o rendezvous com obstáculos físicos por meio de teoria de potenciais virtuais e de um formalismo geométrico baseado em grupos de Lie. Embora esse estudo não trate diretamente da presença de um manipulador robótico, ele contribui para a compreensão de estratégias de guiamento e controle aplicáveis a cenários de aproximação orbital em que a segurança da trajetória e a orientação da espaçonave são aspectos críticos \cite{fiori2024modeling}.

De modo geral, a literatura mostra que o planejamento de trajetória em missões espaciais vem incorporando, de forma crescente, múltiplos objetivos, restrições complexas e diferentes graus de autonomia. Ainda assim, muitos trabalhos tratam separadamente o planejamento do movimento do manipulador, o guiamento da plataforma ou a navegação em ambientes restritos, sem integrar de forma explícita a dinâmica orbital relativa, a atitude da espaçonave e o comportamento acoplado do manipulador em uma única formulação. Essa observação é relevante para a presente dissertação, pois reforça a necessidade de abordagens integradas para as fases de aproximação e interação em operações de berthing \cite{santos2022machine,capolupo2017heuristic,fiori2024modeling,santos2023physics}.